\section{线程粗化与性能权衡}
\label{l13:sec:coarsening}

\subsection{增加每个线程处理的输入数}

线程粗化（thread coarsening）让一个线程承担多个原先可以分配给不同线程的输入元素。在已经使用网格跨步循环的直方图中，可以保持每块线程数 $T$ 不变，减小网格块数 $G$，让每个线程自然执行更多轮循环。每个块仍只建立一份私有直方图，结束时也只合并这一份结果。

设希望每个线程大致处理 $c$ 个元素，一种网格选择为
\begin{equation}
G_c=\left\lceil\frac{N}{cT}\right\rceil,
\qquad P_c=G_cT.
\label{l13:eq:coarsening}
\end{equation}
该公式用于选择工作划分；末尾不足一轮由边界条件处理。只要 $N>0$，线程 $g$ 继续按 $i=g+kP_c$ 遍历即可。每一轮中相邻线程仍读取相邻地址，因此这种粗化保留了输入访问的合并模式。

在固定输入长度下，增加 $c$ 使私有直方图数量由约 $N/T$ 减为约 $N/(cT)$。清零量和不跳过零桶时的归并量从约 $NB/T$ 降为 $NB/(cT)$。若跳过零桶，确切归并量仍是 $\sum_rU_r$，但它的 $G_cB$ 上界随块数减少。

这种收益来自实际减少线程块数量。仅把循环展开、让代码一次写出多次更新，而启动的块数保持不变，并不会减少局部副本数量及每块合并开销。另一方面，循环变长也不会减少总输入读取和总共享原子加：它们仍各为 $N$ 次。

\subsection{粗化的数量例子与并行度代价}

设 $N=2^{20}=1\,048\,576$ 个输入，$B=256$ 个桶，每块 $T=256$ 个线程。表\ref{l13:tab:coarsening-counts}比较 $c=1$ 与 $c=8$ 的网格设置。两种设置都沿用代码\ref{l13:lst:private-hist}，没有改变桶数或局部数组类型。

\begin{table}[htbp]
\centering
\caption{固定输入规模下减少块数的效果}
\label{l13:tab:coarsening-counts}
\small
\begin{tabularx}{\textwidth}{Yrr}
\toprule
数量 & $c=1$ & $c=8$ \\
\midrule
线程块数 $G_c$ & $4096$ & $512$ \\
每线程处理的输入数 & $1$ & $8$ \\
平均每块输入数 & $256$ & $2048$ \\
私有桶清零次数 $G_cB$ & $1\,048\,576$ & $131\,072$ \\
全局合并次数的 $G_cB$ 上界 & $1\,048\,576$ & $131\,072$ \\
共享内存原子加次数 & $1\,048\,576$ & $1\,048\,576$ \\
每块共享直方图容量 & $1\,\mathrm{KiB}$ & $1\,\mathrm{KiB}$ \\
\bottomrule
\end{tabularx}
\end{table}

清零工作量和全局合并上界下降为原来的八分之一，每个块能用同一份共享数组聚合更多贡献。对于全部输入集中于一个桶的极端情况，实际全局归并次数分别为 $4096$ 和 $512$。对于分散输入，实际次数由各块观察到的非零桶数决定，不能把表中的上界直接当作测得次数。

增加粗化程度也会让单个块持续更久。如果所有输入集中于一个桶，每块内部对该桶的串行链随局部输入数增加；若整个网格只剩很少几个块，一些 SM 可能没有足够工作，设备无法利用全部并行资源。因而不能仅凭全局归并次数最少，就断定最快的设置应当只使用一个线程块。

通过延长同一个跨步循环实现粗化时，不必然需要更多寄存器，因为循环中可以反复使用同一组索引和临时变量。若具体实现同时保留多组输入、中间结果或展开后的独立状态，才可能增加寄存器压力。粗化的资源代价需要对应到实际代码，而它减少私有副本数量的机制由网格规模直接决定。

\subsection{线程级局部聚合}

粗化使一个线程依次处理多个输入，也提供了在线程内部合并更新的机会。若为每个线程保存完整直方图，需要每线程 $B$ 个计数器；桶数较大时，寄存器或局部存储开销可能过高。一个更小的状态只保存最近访问的桶编号和该桶连续出现的次数。线程再次读到同一桶时只增加私有计数，读到不同桶时才把上一段计数原子加入块内直方图。

设线程 $g$ 实际读取的桶序列为 $b_{g,0},\ldots,b_{g,m_g-1}$。将其中连续相同的桶划为一段，段数为
\begin{equation}
R_g=\begin{cases}
0,&m_g=0,\\
1+\displaystyle\sum_{k=1}^{m_g-1}\mathbf{1}\{b_{g,k}\ne b_{g,k-1}\},&m_g>0.
\end{cases}
\label{l13:eq:runs}
\end{equation}
每一段只提交一次，因此共享内存原子加的次数由 $N$ 变为 $\sum_g R_g\le N$。每个桶的贡献总量不变，块内直方图及其非零桶数也不变，所以最后的全局合并次数仍为 $\sum_r U_r$。

以下伪代码替换共享直方图核函数中的输入统计循环；前面的清零与第一处屏障，以及后面的第二处屏障与全局合并，均保持不变。\mbox{\code{valid}} 表示是否存在尚未提交的一段，不能用某个合法桶编号表示空状态。

\noindent\begin{minipage}{\textwidth}
\begin{lstlisting}[caption={线程内合并连续同桶更新的伪代码},label={l13:lst:aggregate},language={}]
valid = false
count = 0
for each valid input index assigned to this thread:
    bin = input[index]
    if valid and bin == last_bin:
        count = count + 1
    else:
        if valid:
            atomicAdd(shared_hist[last_bin], count)
        last_bin = bin
        count = 1
        valid = true
if valid:
    atomicAdd(shared_hist[last_bin], count)
// All threads then reach the block barrier.
\end{lstlisting}
\end{minipage}

线程读到的序列若为 $(7,7,2,2,2,7)$，就依次提交 $(7,2)$、$(2,3)$ 和 $(7,1)$ 三个“桶、增量”对，将六次共享原子加减少为三次。最后一对在循环结束后提交；省略这一步会丢失末段。空线程不提交任何计数，但仍参加块内屏障；只处理一个元素的线程则通过末段提交完成唯一一次更新。私有段计数和共享计数器同样必须能表示实际累计数量。

这里的连续性取决于线程的访问顺序。网格跨步循环使同一线程依次读取 $g,g+P,g+2P,\ldots$，原数组中相邻像素的相似性不保证这些位置也相似。例如，输入 $(0,0,1,1,0,0,1,1)$ 由 $P=2$ 个线程跨步读取时，每个线程都看到 $(0,1,0,1)$，八次原子更新一次也不能省略；若改成 $(0,1,0,1,0,1,0,1)$，两个线程分别看到四个 $0$ 和四个 $1$，只需两次提交。输入的桶频数在这两个例子中相同，聚合收益却不同。

当线程所见序列包含较长的同桶段时，少量私有状态能够明显减少热点桶上的共享原子操作。若相邻访问很少命中同桶，$R_g$ 接近 $m_g$，原子操作数量几乎不变，而比较、分支与状态维护仍有成本。因此，线程级聚合应结合实际访问顺序和输入分布评估，不能仅凭直方图中存在高频桶就断定会加速。

\subsection{访存、分支和复用的整体判断}

输入数组通常只被读取一次，每个元素只用于决定自己的桶，因此算法没有需要专门缓存的跨线程输入复用。私有直方图则不同：同一个块内的许多输入可能反复更新同一组桶，输出状态具有复用。将这组状态放进共享内存，正是把访问频繁的数据保留在块内。

输入跨步循环的边界条件主要影响最后一轮不满的线程束，或者没有输入任务的线程。合并阶段的非零判断也可能使一部分线程提交更新、另一部分跳过；这种分支换来了对零贡献原子操作的省略。相比之下，很多线程同时命中同一桶造成的是原子地址竞争：即使所有线程执行完全相同的指令，也可能因为目标地址相同而排队。

\begin{table}[htbp]
\centering
\caption{私有化直方图的主要性能因素}
\label{l13:tab:performance-factors}
\small
\begin{tabularx}{\textwidth}{>{\raggedright\arraybackslash}p{2.0cm}YY}
\toprule
观察对象 & 形成的开销 & 对应调整 \\
\midrule
同轮输入地址 & 地址分散时需要更多内存事务 & 使相邻线程读取相邻输入 \\
公共桶竞争 & 多块贡献排成同地址更新序列 & 先在块内聚合，再提交局部计数 \\
块内热点桶 & 共享内存同址原子更新仍排队 & 结合局部输入量与数据分布评估 \\
副本数量 & 每块都需清零和逐桶合并 & 适度增加每块工作量、减少块数 \\
驻留与并行度 & 资源需求大或网格过小会限制并行 & 同时考虑桶容量、线程数和块数 \\
\bottomrule
\end{tabularx}
\end{table}

直接全局原子、共享内存私有化和线程粗化分别改变不同环节。基本算法提供简单的正确实现；私有化改变竞争发生的位置及公共更新次数；粗化进一步摊薄每份私有副本的固定开销。比较这些方案时，应固定输入长度、桶映射和数据分布，再结合操作数量与实际运行时间判断收益。
